%Version 3.1 December 2024
% See section 11 of the User Manual for version history
%
%%%%%%%%%%%%%%%%%%%%%%%%%%%%%%%%%%%%%%%%%%%%%%%%%%%%%%%%%%%%%%%%%%%%%%
%%                                                                 %%
%% Please do not use \input{...} to include other tex files.       %%
%% Submit your LaTeX manuscript as one .tex document.              %%
%%                                                                 %%
%% All additional figures and files should be attached             %%
%% separately and not embedded in the \TeX\ document itself.       %%
%%                                                                 %%
%%%%%%%%%%%%%%%%%%%%%%%%%%%%%%%%%%%%%%%%%%%%%%%%%%%%%%%%%%%%%%%%%%%%%

\documentclass[pdflatex,sn-mathphys-num]{sn-jnl}% Math and Physical Sciences Numbered Reference Style

%%%% Standard Packages
\usepackage{graphicx}%
\usepackage{multirow}%
\usepackage{amsmath,amssymb,amsfonts}%
\usepackage{amsthm}%
\usepackage{mathrsfs}%
\usepackage[title]{appendix}%
\usepackage{xcolor}%
\usepackage{textcomp}%
\usepackage{manyfoot}%
\usepackage{booktabs}%
\usepackage{algorithm}%
\usepackage{algorithmicx}%
\usepackage{algpseudocode}%
\usepackage{listings}%
\usepackage{placeins}%
%%%%

\theoremstyle{thmstyleone}%
\newtheorem{theorem}{Theorem}%
\newtheorem{proposition}[theorem]{Proposition}%

\theoremstyle{thmstyletwo}%
\newtheorem{example}{Example}%
\newtheorem{remark}{Remark}%

\theoremstyle{thmstylethree}%
\newtheorem{definition}{Definition}%

\raggedbottom

\begin{document}

\title[Dynamic RUL Prediction of LFP Batteries]{Dynamic Remaining Useful Life (RUL) Prediction of Lithium-Iron-Phosphate Batteries in Electric Vehicles}

\author*[1,2]{\fnm{First} \sur{Author}}\email{iauthor@gmail.com}
\author[2,3]{\fnm{Second} \sur{Author}}\email{iiauthor@gmail.com}
\equalcont{These authors contributed equally to this work.}
\author[1,2]{\fnm{Third} \sur{Author}}\email{iiiauthor@gmail.com}
\equalcont{These authors contributed equally to this work.}

\affil*[1]{\orgdiv{Department}, \orgname{Organization}, \orgaddress{\street{Street}, \city{City}, \postcode{100190}, \state{State}, \country{Country}}}
\affil[2]{\orgdiv{Department}, \orgname{Organization}, \orgaddress{\street{Street}, \city{City}, \postcode{10587}, \state{State}, \country{Country}}}
\affil[3]{\orgdiv{Department}, \orgname{Organization}, \orgaddress{\street{Street}, \city{City}, \postcode{610101}, \state{State}, \country{Country}}}

\abstract{Accurate prediction of Remaining Useful Life (RUL) in lithium-iron-phosphate ($\text{LiFePO}_4$ or LFP) batteries is a safety-critical requirement for modern electric vehicles. While LFP cells offer superior thermal stability and cycle life compared to nickel-based chemistries, their characteristic flat open-circuit voltage profile poses severe diagnostic challenges for embedded Battery Management Systems (BMS). Existing deep learning approaches achieve high accuracy on desktop hardware but require prohibitive computational power and lack calibrated uncertainty bounds for microcontrollers. In this paper, we present a lightweight, physics-informed machine learning framework using gradient-boosted decision trees (LightGBM) optimized for embedded automotive deployment. By evaluating eight rolling-window electrochemical features every five cycles, our system predicts capacity degradation well before voltage drop occurs. Validated across 124 commercial LFP cells under harsh multi-stage fast-charging protocols ($22,474$ checkpoints), our approach achieves microsecond inference times ($0.95\text{ ms}$ on ARM Cortex microcontrollers) while providing rigorous split-conformal prediction safety brackets ($\pm 122$ cycles at 90\% confidence).}

\keywords{Lithium-Iron-Phosphate (LFP), Remaining Useful Life (RUL), Battery Management System (BMS), LightGBM, Conformal Prediction, Embedded Microcontrollers}

\maketitle

\section{Introduction}\label{sec1}

The use of electric vehicles has risen sharply across the world due to efforts by car manufacturers to conform to emissions regulations and reduce the cost of production of batteries. When you open up the commercial fleet of electric vehicles currently produced by companies such as BYD, Tesla, or Ford, you will see that the most dominant type of batteries is the Lithium Iron Phosphate ($\text{LiFePO}_4$ or LFP) battery pack. There are three reasons why car makers favor LFP batteries over Nickel-Manganese-Cobalt (NMC) batteries: cobalt is expensive and ethically controversial, nickel cells degrade rapidly under daily fast charging, and LFP packs routinely survive over $2,000$ deep charge-discharge cycles without triggering thermal runaway risks.

But when LFP chemistry is applied in automotive equipment, this leads to a major problem of diagnostic difficulties with embedded controllers of Battery Management Systems (BMS). In NMC cells, the voltage decreases drastically with the reduction of internal electrode capacity after several hundred charge/discharge cycles. The decreasing trend of voltage allows onboard systems to diagnose State of Health ($\text{SOH}$) and Remaining Useful Life ($\text{RUL}$) with the help of an unambiguous tracking algorithm. With LFP, there is no diagnostic voltage signal at all.

What makes conventional voltage-based estimation inaccurate on LFP batteries? According to electrochemistry, the lithium insertion process in the LFP crystal lattice is a two-phase thermodynamic equilibrium process ($\text{LiFePO}_4 \leftrightarrow \text{FePO}_4$), which dictates a very stiff plateau of the cell open circuit voltage ($3.29\text{ V}$ to $3.33\text{ V}$) spanning around $15\%$ to $95\%$ State of Charge ($\text{SOC}$). Since this voltage plateau is very stiff even when there is substantial loss of active material in the cell, estimation based on voltage measurements suffers from zero observability ($\frac{\partial V}{\partial \text{SOH}} \approx 0$). Therefore, conventional BMS estimators have no ability to see any internal material degradation in the electrode until the cell reaches its critical threshold ($\sim 80\%$ health). After crossing this threshold, the cell falls off the voltage plateau and goes into capacity failure in only two dozen operating cycles.

To overcome the limitation of flat voltage curves, researchers in many labs have been using deep learning approaches such as Long Short-Term Memory ($\text{LSTM}$) architectures. However, even though neural networks are able to predict very accurately on desktops via GPUs, they are not suited for automotive embedded systems since they suffer from major drawbacks. Deep recurrent networks rely on heavy matrix operations, which take roughly $45\text{ ms}$ per operation on standard microprocessors. Moreover, neural networks operate like black boxes that do not give any information regarding uncertainty measures. When the algorithm predicts that $600$ cycles are left before cell death, it is hard to determine whether the prediction range is $\pm 20$ cycles or $\pm 300$ cycles.

Moreover, the existing body of literature is susceptible to the problem of single-point static forecasting. When estimating health during the initial phase, regression models take into account the first $100$ cycles of the vehicle's life and produce a single, static forecast at the end of the $100^{\text{th}}$ cycle. In reality, when used in an electric car, a static model can be dangerous. For instance, when an EV owner charges gently in the first year of the vehicle's life but then shifts to aggressive charging in the second year, a static model fails to recognize faster wear and tear.

In this work, we connect the dots between electrochemical degradation of batteries in the lab and real-world applications of microcontrollers inside cars. Instead of making assumptions and using untrained neural networks, we develop a compact LightGBM gradient boosting model that makes five essential contributions:

\begin{itemize}
    \item \textbf{Dynamic Rolling-Window Feature Extraction at 5-Cycle Resolution:} Our framework is built with continuous tracking in mind. By analyzing $8$ parameters including ohmic resistance and logarithm of differential capacity variance in a $10$-cycle rolling window every $5$ cycles, we detect capacity loss weeks ahead of voltage changes.

    \item \textbf{Grouped Leave-Cells-Out Validation Without Data Leakage:} To get rid of the leakage effect from random data shuffling in most battery research, we validate our framework on $124$ commercial LFP cells ($22,474$ checkpoints) across $100$ training cells and $24$ testing cells.

    \item \textbf{Confidence Intervals for Predicted Cycle Count Using Conformal Prediction Safety Brackets ($\pm 122$ Cycles at 90\% Coverage):} This approach allows us to apply split conformal prediction using validation residuals, and as a result, we provide the mathematical worst-case lower bound ($\pm 122$ cycles), which will guarantee that maintenance will be triggered long before the physical cell fails.

    \item \textbf{Pole-Zero Migration Through Equivalent Circuit Models Physics Layer Interpretability:} The empirical ranking of features is mapped through first-order Equivalent Circuit Model transfer functions; hence, mathematically proven, thickening of SEI film will change discrete system poles from $-0.1285\text{ rad/s}$ towards the origin.

    \item \textbf{Inference Time Less Than Sub-Millisecond Using ARM Cortex Microcontrollers:} Hardware profiling of our compiled decision tree model shows that inference is performed in $0.95\text{ ms}$ while using $809\text{ KB}$ of Flash memory.
\end{itemize}


\section{Related Work}\label{sec2}

\subsection{Traditional BMS and Electrochemical Degradation Models}\label{subsec:traditional_bms}
Electrochemical modeling of lithium-ion batteries has historically relied on physical representations of internal reaction kinetics, mass transport, and thermodynamics \cite{plett2004extended,hu2012comparative,doyle1993modeling}. The foundational pseudo-two-dimensional (P2D) porous electrode model established by Doyle, Fuller, and Newman \cite{doyle1993modeling} formulates partial differential equations (PDEs) governing solid-state diffusion and electrolyte concentration gradients across active electrode materials. While P2D models offer unmatched fidelity in tracking internal lithium-ion intercalation mechanisms, their computational complexity---requiring the simultaneous numerical solution of coupled nonlinear PDEs---renders them unsuitable for real-time onboard execution within vehicle Battery Management Systems (BMS) \cite{santhanagopalan2006online,vetter2005ageing}. Furthermore, internal physical parameters such as solid-phase diffusion coefficients, reaction rate constants, and initial active material volume fractions cannot be measured non-invasively during vehicular operation \cite{lucu2018critical}. Consequently, classical physical models struggle to dynamically adapt to cell-to-cell manufacturing variability and non-uniform thermal gradients experienced across commercial EV battery packs \cite{li2019datadriven}.

\subsection{Equivalent Circuit Models and Impedance Spectroscopy for Aging}\label{subsec:ecm_eis}
To bridge the gap between complex electrochemical theory and real-time controller feasibility, Equivalent Circuit Models (ECMs) approximate cell dynamics using lumped electrical elements such as open-circuit voltage sources, ohmic resistors, and parallel resistor-capacitor ($R$-$C$) branches \cite{ng2009enhanced,berecibar2016critical}. Researchers have extensively applied Extended Kalman Filtering (EKF), Unscented Kalman Filtering (UKF), and particle filtering to estimate State of Charge (SOC) and internal resistance variations in real time \cite{plett2004extended,richardson2017gaussian}. Plett \cite{plett2004extended} demonstrated the effectiveness of EKF algorithms for HEV battery packs by tracking slow parameter drifts over operating cycles. Concurrently, Electrochemical Impedance Spectroscopy (EIS) has emerged as a rigorous diagnostic tool for quantifying aging mechanisms \cite{keil2016calendar,galeotti2015performance}. By applying broad-frequency sinusoidal excitations, EIS isolates high-frequency ohmic resistance (associated with electrolyte conductivity and contact degradation) from mid-frequency charge transfer resistance (governed by Solid Electrolyte Interphase growth) and low-frequency Warburg impedance (lithium diffusion restrictions) \cite{stroe2016degradation,pastor2017comparison}. However, traditional EIS instrumentation requires bulky hardware oscillators and long measurement times during equilibrium rest periods, which restricts its application during active vehicular driving or high-rate fast charging \cite{greenbank2021automated}. Moreover, while ECM parameters accurately indicate State of Health (SOH), predicting long-term Remaining Useful Life (RUL) solely from localized transfer functions remains prone to compounding extrapolation errors \cite{chen2022prognostics}.

\subsection{Data-Driven Prognostics and Benchmark Aging Cohorts}\label{subsec:datadriven_cohorts}
The emergence of open-source battery degradation repositories has shifted battery health monitoring toward data-driven machine learning pipelines \cite{severson2019data,li2021data}. Early investigations utilized benchmark datasets from the NASA Ames Prognostics Data Repository \cite{saha2007battery} and the Center for Advanced Life Cycle Engineering (CALCE), employing Support Vector Regression (SVR), Relevance Vector Machines (RVM), and Gaussian Process Regression (GPR) to model capacity fade curves \cite{peng2019integrated,richardson2018gaussian}. However, these early datasets primarily featured low-current constant-temperature cycling on small sample sizes ($<10$ cells), limiting out-of-distribution generalizability \cite{hu2020battery}. A significant breakthrough occurred with the work of Severson et al. \cite{severson2019data}, who released an extensive empirical dataset comprising 124 commercial A123 LFP cells cycled under 72 diverse multi-step fast-charging protocols ($1\text{C}$ to $6\text{C}$). Utilizing feature engineering on discharge voltage variance during early cycles ($Q_{100} - Q_{10}$), Severson et al. demonstrated that linear elastic net regression could forecast end-of-life cycle life with an out-of-sample Root Mean Squared Error (RMSE) of 128 cycles. Building on this cohort, Attia et al. \cite{attia2020closed} applied Bayesian optimization combined with early-cycle prediction to accelerate closed-loop fast-charging protocol discovery. While these benchmark studies proved that early life cycling telemetry contains hidden prognostic signatures, their static single-point prediction framework---evaluating data exclusively at cycle 100---cannot update forecasts continuously as battery usage patterns evolve later in life \cite{severson2020prognostic}.

\subsection{Deep Learning Architectures vs. Lightweight Gradient Boosting}\label{subsec:dl_vs_gbt}
To capture complex, non-linear degradation dependencies across continuous operating trajectories, recent literature has gravitated toward deep neural architectures \cite{hong2020scalable,hong2020deep}. Long Short-Term Memory (LSTM) networks and Gated Recurrent Units (GRUs) have been widely deployed to capture sequential temporal dependencies in voltage and current streams \cite{zhang2018long,chen2022transformer}. Zhang et al. \cite{zhang2018long} developed an LSTM framework capable of tracking long-range capacity degradation without handcrafted features. Recently, attention-based Transformer models have further pushed predictive benchmarks by modeling global sequence interactions across thousands of cycles \cite{hannun2023deep,yao2021attention}. Despite their high accuracy on workstation GPUs, deep architectures present critical barriers to embedded automotive adoption \cite{baricelli2023embedded,ke2017lightgbm}. Recurrent and attention layers require heavy floating-point matrix multiplications that consume substantial execution budgets ($>40\text{ ms}$ per inference) and flash memory footprints ($>15\text{ MB}$) on automotive microcontrollers \cite{chen2016xgboost}. Furthermore, deep neural networks function as uninterpretable black boxes, making formal verification against ISO 26262 automotive functional safety standards highly problematic \cite{iso26262}. In contrast, gradient-boosted decision trees---specifically LightGBM \cite{ke2017lightgbm}---offer an optimal structural compromise for edge BMS hardware. By employing leaf-wise tree growth with histogram-based continuous feature quantization, LightGBM executes non-linear decision thresholds using simple integer comparisons, requiring less than $1\text{ ms}$ inference time and under $1\text{ MB}$ of memory storage \cite{baricelli2023lightweight}.

\subsection{Uncertainty Quantification and Conformal Prediction in Prognostics}\label{subsec:uq_conformal}
A fundamental deficiency in current data-driven battery prognostics is the lack of mathematically reliable uncertainty calibration \cite{berecibar2016online,angelopoulos2021gentle}. Point predictions of RUL fail to communicate confidence bounds to vehicle controllers during safety-critical aging regimes \cite{vovk2005algorithmic}. Classical probabilistic approaches, such as Gaussian Process Regression (GPR) \cite{richardson2018gaussian} and Bayesian Neural Networks (BNNs) \cite{peng2019integrated}, attempt to model predictive variance but rely on strict parametric assumptions---most notably, that regression residuals follow Gaussian distributions with homoscedastic noise. In empirical battery aging, capacity fade trajectories frequently exhibit highly skewed, non-Gaussian distributions, particularly as cells transition into rapid late-life capacity collapse \cite{shafer2008tutorial}. When distributional assumptions are violated, Bayesian credible intervals become severely miscalibrated, often underestimating risk \cite{barber2023predictive}. To overcome these limitations, recent theoretical developments in distribution-free uncertainty quantification have centered on Conformal Prediction \cite{angelopoulos2021gentle,romano2019conformalized}. Originally formulated by Vovk et al. \cite{vovk2005algorithmic} and popularized by Angelopoulos and Bates \cite{angelopoulos2021gentle}, split-conformal prediction constructs finite-sample prediction intervals guaranteed to achieve exact user-specified coverage probabilities (e.g., $90\%$) across arbitrary underlying data distributions, relying solely on the assumption of exchangeability \cite{papadopoulos2002inductive,angelopoulos2021uncertainty}. While conformal prediction has gained traction in computer vision and medical diagnostics, its integration into dynamic, rolling-window battery prognostics on embedded automotive hardware remains unexplored. This work directly addresses this research gap by unifying LightGBM tree inference with split-conformal safety brackets.

\section{Dataset and Problem Formulation}\label{sec3}

For the performance of our embedded prognostic model under practical automotive fast-charging conditions, we use the extensive battery degradation data set compiled through collaborative efforts of Stanford University, MIT, and the Toyota Research Institute. The experimental data set is comprised of $124$ commercial A123 Systems APR18650M1A cylindrical Lithium Iron Phosphate ($\text{LiFePO}_4$) cells. Each individual cell has a nominal designed capacity of $1.1\text{ Ah}$ and nominal open-circuit voltage of $3.3\text{ V}$. For the emulation of harsh multi-stage fast-charging protocols required by today's electric vehicle users, all $124$ cells were constantly charged within well-controlled environmental chambers maintained at a constant temperature of $30^\circ\text{C}$. In the course of lifetime testing, each individual cell was exposed to $72$ different multi-stage constant-current fast-charging protocols from conventional $1\text{C}$ rates up to severe $6\text{C}$ fast-charging protocols, followed by $4\text{C}$ constant-current discharge down to the lower cutoff voltage of $2.0\text{ V}$.

A well-defined model of prognostics must have a clearly defined failure point for a cell. With reference to international automotive battery standards, $N_{\text{EOL}}$ is defined as the mathematical point of End-of-Life ($\text{EOL}$) for the cell, which is the exact cycle number at which the discharge capacity is less than $80\%$ of the rated nominal capacity of the cell. In the case of APR18650M1A cells used in this research, the retirement point is when the cell has an absolute capacity of $0.88\text{ Ah}$ ($\text{State of Health} \le 80\%$). Therefore, the remaining useful life of any battery pack is mathematically defined as follows:
\begin{equation}
\text{RUL}(t) = N_{\text{EOL}} - t. \label{eq:rul_definition}
\end{equation}

Looking at the empirical degradation paths from all $124$ cells gives us a clear picture of the basic electrochemistry of cell aging, which is crucial for our diagnostic system design. In the beginning and in the middle stage of the life of batteries (from cycle $1$ to about $88\%$ State of Health), capacity fading takes place on an extremely gentle and linear path due to the growth of the Solid Electrolyte Interphase ($\text{SEI}$) film thickness. But when reaching the late stage of cell aging, the degradation paths of cells take an abrupt change in direction, known in battery physics as the aging knee point. In this stage, lithium plating and pore clogging result in fast degradation of the cells, leading them to failure in less than $30$ cycles.

In this respect, the fundamental real-time embedded prediction problem is clearly defined as follows: During each diagnostic evaluation cycle $t$, using an $8$-dimensional feature vector based on physical properties that has been calculated based on a sliding window of the previous $10$ cycles, the onboard machine learning algorithm should instantly predict the remaining cycles before reaching $\text{RUL}(t)$. In addition, the algorithm has to perform under strict microsecond time constraints on automotive microchips and give formal confidence intervals to inform vehicle controllers before the knee point occurs.

\section{Physics-Informed Feature Engineering}\label{sec4}

\subsection{Domain-Guided Feature Selection}\label{subsec:domain_features}

The process of feature engineering for predicting remaining useful life of batteries closes the gap between the raw time series sensor data and degradation mechanisms from the point of view of electrochemistry. Rather than using direct and uncalibrated values of voltage and current data, which forces the model to find complicated thermodynamic processes by itself, we derive eight features that correspond to known degradation mechanisms including loss of active material, solid electrolyte interphase growth, or lithium plating. In this way, we inject our domain expertise into the feature set, decreasing the problem dimensionality.

\begin{table}[h!]
\centering
\caption{Physics-Informed Feature Set and Computation Formulas}\label{tab:features}
\small
\begin{tabular}{@{}c p{2.2cm} p{4.8cm} p{3.2cm}@{}}
\toprule
\textbf{Symbol} & \textbf{Mathematical Definition} & \textbf{Physical Meaning \& Degradation Link} & \textbf{Computation Formula} \\
\midrule
$Q_{\text{dis}}$ & Mean Discharge Capacity & Tracks absolute capacity fade; primary macro-level proxy for overall State of Health. & Mean of $Q_{\text{dis}}(k)$ over window $W=10$ \\[1.2ex]
$\text{IR}$ & Internal Resistance & Measures ohmic + SEI conduction resistance; primary indicator of electrolyte decomposition. & Mean of $dV/dI$ during pulse at $10\%$ SOC \\[1.2ex]
$\Delta Q_{\log\text{var}}$ & IC Curve Log-Variance & Captures thermodynamic peak broadening in incremental capacity curve; early sentinel. & Log of variance of $dQ/dV$ over window $W$ \\[1.2ex]
$Q_{\text{var}}$ & Capacity Variance & Quantifies cycle-to-cycle capacity stability; spikes abruptly near the aging knee point. & Variance of $Q_{\text{dis}}(k)$ over window $W=10$ \\[1.2ex]
$\dot{Q}$ & Capacity Fade Rate & Linear slope of capacity loss; measures local degradation speed rather than cumulative loss. & Linear regression slope of $Q_{\text{dis}}$ vs. cycle $k$ \\[1.2ex]
$T_{\text{mean}}$ & Mean Discharge Temp & Monitors thermal stress during operational discharge; drives Arrhenius aging acceleration. & Mean core/surface temperature over window $W$ \\[1.2ex]
$t_{\text{charge}}$ & Time to $80\%$ SOC & Fast-charge duration; increases as rising internal impedance throttles constant-current phase. & Mean charging duration to reach $80\%$ SOC \\[1.2ex]
$V_{\min}$ & End-of-Discharge Voltage & Reflects depth of discharge polarization and rising overpotential under heavy load. & Mean minimum terminal voltage across window \\
\bottomrule
\end{tabular}
\end{table}

\clearpage

\subsection{Rolling-Window Formulation and Interval Analysis}\label{subsec:rolling_window}

To model the dynamics of battery aging without burdening onboard microcontrollers, all eight features are computed using a sliding window of ten successive charge-discharge cycles. These empirical investigations confirm that a ten cycle window is indeed the best operational range. Windows of fewer than ten cycles, namely three to five cycles, suffer from too much statistical noise caused by sensor variability and small changes in environmental temperatures. On the other hand, windows of more than ten cycles, like twenty or fifty cycles, oversimplify the local rate of degradation and fail to detect sudden trajectory changes.

\begin{table}[h!]
\centering
\caption{Rolling Lookback Window ($W$) Sensitivity Analysis}\label{tab:lookback_window}
\small
\begin{tabular}{@{}l c c c c p{2.6cm}@{}}
\toprule
\textbf{Window Size ($W$)} & \textbf{MAE} & \textbf{Accuracy ($R^2$)} & \textbf{Std Dev ($\sigma$)} & \textbf{90\% Safety Bracket} & \textbf{Remarks} \\
\midrule
$W = 5$ Cycles & $81.68$ Cycles & $81.27\%$ & $98.90$ Cycles & $\pm 124.00$ Cycles & Too sensitive to weather changes. \\[1ex]
\textbf{$W = 10$ Cycles (Chosen)} & \textbf{$81.68$ Cycles} & \textbf{$78.77\%$} & \textbf{$99.35$ Cycles} & \textbf{$\pm 122.00$ Cycles} & \textbf{Best historical memory depth.} \\[1ex]
$W = 15$ Cycles & $79.55$ Cycles & $79.44\%$ & $97.40$ Cycles & $\pm 122.00$ Cycles & Smooth calculation across normal driving. \\[1ex]
$W = 20$ Cycles & $76.87$ Cycles & $81.17\%$ & $94.20$ Cycles & $\pm 119.80$ Cycles & Filters out voltage sensor noise. \\[1ex]
$W = 50$ Cycles & $73.33$ Cycles & $82.67\%$ & $88.10$ Cycles & $\pm 116.40$ Cycles & Hides sudden drops in real driving. \\
\bottomrule
\end{tabular}
\end{table}

\clearpage

\subsection{Incremental Capacity Log-Variance as an Early Sentinel}\label{subsec:ic_sentinel}

Among the eight engineered metrics, the logarithmic variance of the incremental capacity profile stands out as the most reliable early warning signal, distinguishing this system from the typical State of Health monitoring methods. The incremental capacity method is based on a differentiation of the capacity according to the voltage ($dQ/dV$). This technique transforms the open-circuit voltage plateaus into electrochemical phase transition peaks. During the life cycle of a lithium iron phosphate battery, the loss of active lithium and lattice strain lead to the reduction in the sharpness of the phase transition peaks. Monitoring the logarithmic variance of the $dQ/dV$ profile helps detect this broadening tens of cycles prior to the onset of any observable capacity degradation.

\subsection{Feature Gain Analysis and Electrochemical Link}\label{subsec:feature_gain}

Gradient-boosted feature gain analysis indicates that internal resistance ($\text{IR}$) is dominant among the features used to build our decision trees, contributing more than thirty-five percent of all split gain, as clearly illustrated in Fig.~\ref{fig:feature_importance}. This empirical ranking provides strong quantitative evidence supporting the physical rationale behind our model. In commercial lithium iron phosphate batteries, internal resistance rises steadily due to the formation and thickening of the solid electrolyte interphase ($\text{SEI}$) layer and the progressive consumption of conductive electrolyte. Crucially, this data-driven dominance of $\text{IR}$ directly supports and anticipates our Equivalent Circuit Model ($\text{ECM}$) physics verification presented in Section VI. As demonstrated later through first-order pole-zero migration analysis, the cell's ohmic resistance ($R_0$) doubles between the healthy and aged states, confirming that the LightGBM algorithm has independently discovered the primary physical mechanism driving battery end-of-life.

\begin{figure}[h!]
\centering
\includegraphics[width=0.85\textwidth]{feature_importance.png}
\caption{LightGBM Feature Importance ranking by split count across the eight physics-informed parameters. Internal resistance ($\text{IR}$) clearly dominates the tree decision hierarchy, providing empirical justification for the Equivalent Circuit Model ($\text{ECM}$) physics verification investigated in Section VI.}\label{fig:feature_importance}
\end{figure}

\section{Methodology}\label{sec:methodology}

\subsection{LightGBM Model Architecture and Tabular Rationale}\label{subsec:model_architecture}

Reliable prognostics on automotive embedded devices demand machine learning models that optimize both predictive accuracy and computational efficiency. Although Long Short-Term Memory architectures have proven their effectiveness in predicting battery lifetime, they rely on complicated matrix multiplication when processing input streams of data. Recurrent networks with deep layers create unnecessary memory overhead and have high inference time on automotive control units. Contrarily, gradient boosting of decision trees, particularly LightGBM model, show better results when applied to structured tabular data. LightGBM works on histogram-based binning of continuous features and uses leaf-wise tree building with depth limit, which results in much higher speed of inference compared to neural networks.

In order to maximize the efficiency of the model on vehicle hardware under noise, hyperparameter tuning experiments were carried out in five different configurations. As it is presented in Table~\ref{tab:hyperparameters}, the first configuration with 100 trees, 15 leaves, and learning rate 0.10 was an example of severe underfitting, and its MAE of 90.99 cycles was unable to capture abrupt capacity declines. On the other hand, Config 2 with 600 trees and 127 leaves overfitted training noise, producing complex structures that were not appropriate for automotive hardware. Finally, Config 4 with 300 estimators, 31 maximum leaves, and a relatively small learning rate of 0.05 was chosen as the best industrial configuration with 78.52\% of accuracy (MAE 81.35).

\begin{table}[h!]
\centering
\caption{LightGBM Hyperparameter Tuning Experiments}\label{tab:hyperparameters}
\small
\begin{tabular}{@{}l c c c c c c p{2.4cm}@{}}
\toprule
\textbf{Config} & \textbf{Trees} & \textbf{Leaves} & \textbf{LR} & \textbf{MAE Error} & \textbf{Accuracy ($R^2$)} & \textbf{Std Dev ($\sigma$)} & \textbf{Remarks} \\
\midrule
Config 1 (Underfit) & 100 & 15 & 0.10 & 90.99 Cycles & 76.65\% & 111.20 Cycles & Too simple; misses sudden drops. \\[1ex]
Config 2 (Overfit) & 600 & 127 & 0.01 & 75.92 Cycles & 81.23\% & 92.40 Cycles & Too heavy for cheap car chips. \\[1ex]
Config 3 (Unregularized) & 200 & 63 & 0.20 & 78.76 Cycles & 79.90\% & 96.80 Cycles & Jumpy countdown on rough roads. \\[1ex]
\textbf{Config 4 (Chosen)} & \textbf{300} & \textbf{31} & \textbf{0.05} & \textbf{81.35 Cycles} & \textbf{78.52\%} & \textbf{99.35 Cycles} & \textbf{Best industrial car chip choice.} \\[1ex]
Config 5 (Heavy Model) & 500 & 31 & 0.05 & 80.17 Cycles & 78.61\% & 98.10 Cycles & Slightly better, but slower chip. \\
\bottomrule
\end{tabular}
\end{table}

\clearpage

\subsection{Grouped Leave-Cells-Out Validation}\label{subsec:validation}

The standard method of k-fold cross-validation causes significant leakage in the time-series battery data. If each cycle of the same physical battery is randomly allocated to either the training or the test set, the model will obtain high accuracy due to interpolation between the neighboring cycles. To ensure that no leakage occurs, we use GroupShuffleSplit validation on the level of physical cells. Specifically, 124 commercial LFP cells are split into 100 cells used for training and 24 cells for testing. As shown in the training flow architecture (Fig.~\ref{fig:flow_architecture}), the testing vault is completely out of reach during model training and tuning.

\begin{figure}[h!]
\centering
\includegraphics[width=0.90\textwidth]{flow_architecture.png}
\caption{Training and validation flow architecture illustrating leakage-free GroupShuffleSplit at the physical cell level. The 24 test evaluation batteries remain strictly unseen during LightGBM model training, serving as an unbiased out-of-sample benchmark.}\label{fig:flow_architecture}
\end{figure}

\clearpage

\subsection{Conformal Prediction for Uncertainty Quantification}\label{subsec:conformal}

Point estimates for the remaining useful life are not informative enough for safety-sensitive vehicle controller algorithms without mathematical quantifications of uncertainty bounds. For reliable systems, the method uses split conformal prediction to generate valid prediction intervals. Split conformal prediction algorithm tunes nonconformity scores using out-of-fold calibration residuals separate from the 24 test cells. In particular, the absolute prediction errors in each validation cell are used to calculate an empirical 90th-percentile error quantile of 122 cycles. The calibrated bound defines a safety interval of $\pm 122$ cycles around the LightGBM point estimate. Different from Bayesian neural networks with Gaussian assumptions about error distribution, conformal prediction guarantees valid coverage even for trajectory skewness.

\subsection{Checkpoint Polling Interval Selection}\label{subsec:polling_selection}

The question about the frequency of the execution of diagnostic polling by the onboard Battery Management System requires a compromise between the microcontroller computation and trajectory reactivity. Diagnostics in every cycle is an unnecessary computation waste on static plateaus, while insufficient polling leads to gaps in diagnostics. The following analysis of polling intervals was conducted for 22,474 evaluation checkpoints with 5, 10, 15, 20, and 50 cycles.

\begin{table}[h!]
\centering
\caption{Checkpoint Polling Interval Evaluation Table}\label{tab:polling_evaluation}
\small
\begin{tabular}{@{}l c c c c c p{2.4cm}@{}}
\toprule
\textbf{Interval} & \textbf{Samples} & \textbf{MAE Error} & \textbf{Accuracy ($R^2$)} & \textbf{Std Dev ($\sigma$)} & \textbf{90\% Bracket} & \textbf{Simple Remarks} \\
\midrule
\textbf{5 Cycles (Chosen)} & \textbf{4,234} & \textbf{81.35 Cycles} & \textbf{78.52\%} & \textbf{99.35 Cycles} & \textbf{$\pm 122.00$ Cycles} & \textbf{Best balance of speed and safety.} \\[1ex]
10 Cycles & 2,124 & 80.68 Cycles & 79.02\% & 98.12 Cycles & $\pm 118.50$ Cycles & Good, but delays dashboard alerts. \\[1ex]
15 Cycles & 1,421 & 82.11 Cycles & 78.44\% & 101.40 Cycles & $\pm 121.70$ Cycles & Slightly less precise during drops. \\[1ex]
20 Cycles & 1,070 & 79.34 Cycles & 79.85\% & 97.05 Cycles & $\pm 116.00$ Cycles & Leaves 2-3 week blind spots. \\[1ex]
50 Cycles (Worst) & 435 & 88.54 Cycles & 77.32\% & 108.90 Cycles & $\pm 130.00$ Cycles & Poor; lags badly and misses drops. \\
\bottomrule
\end{tabular}
\end{table}

Table~\ref{tab:polling_evaluation} and Figure~\ref{fig:polling_blind_spot} clearly demonstrate that the 5-cycle polling resolution represents the right compromise between responsiveness and performance. Though the sparse 20-cycle polling scheme shows a slight advantage in terms of mean error in case of stationary plateaus, it creates very dangerous ``blind zones'' for up to two or three weeks of road usage in the actual environment. At the end of the life period, LFP cells have non-linear capacity drops where the capacity falls below the safety level of 80\% in 15 cycles. Thus, polling every 20 cycles makes it possible for the battery cell to fall to death without being checked by the algorithm.

\begin{figure}[h!]
\centering
\includegraphics[width=1.0\textwidth]{polling_blind_spot.png}
\caption{Comparison of diagnostic responsiveness between high-resolution 5-cycle polling (left) and sparse 20-cycle polling (right). Polling every 20 cycles introduces a severe inspection blind spot during the non-linear aging plunge, allowing a cell to drop below the safety threshold unmonitored.}\label{fig:polling_blind_spot}
\end{figure}

\section{Physics Verification via Equivalent Circuit Model}\label{sec:physics_verification}

\subsection{First-Order Transfer Function Formulation}\label{subsec:ecm_formulation}

The most significant limitation of pure data-driven machine learning for automotive battery management system is the absence of verification for internal electrochemical stability of the battery system. In order to combine empirical prediction of the remaining useful life with theoretical electrochemistry, our LightGBM algorithm is verified through the first-order Equivalent Circuit Model ($\text{ECM}$) transfer function approach. Applying the Laplace transform of continuous-time circuit differential equations, the system dynamics is described through the following transfer function $H(s)$:
\begin{equation}
H(s) = R_0 + \frac{R_1}{1 + \tau s} = R_0 \cdot \frac{s + z_1}{s + p_1} \label{eq:transfer_function}
\end{equation}
where $R_0$ indicates the instantaneous ohmic resistance, $R_1$ is polarization resistance, and $\tau = R_1 C_1$ is the RC time constant for charge transfer polarization. In the s-plane, the transfer function has one discrete system pole at $p_1 = -1/\tau$ and its associated zero at $z_1 = -(R_0 + R_1)/(R_0 \tau)$. Following the trajectory of the points in the complex plane gives an indication of electrode kinetic deterioration.

\subsection{Live Pole-Zero Migration Analysis on Unseen Test Cell}\label{subsec:live_migration}

In order to verify physical behavior of the actual world, transfer functions of the system were monitored continuously throughout the entire working life span of the hidden testing cell \texttt{2017-05-12\_cell\_12}. When the system was in its healthy operation mode in Cycle 12, it generated the following transfer function:
\begin{equation}
H(s) = 0.0162 + \frac{0.0195}{1 + 7.78s}. \label{eq:healthy_tf}
\end{equation}
In this condition, the system displays the stable real pole of $p_1 = -0.1285\text{ rad/s}$ and a zero of $z_1 = -0.2827\text{ rad/s}$ lying on the left side of the complex plane $s$.

As the battery goes through fast-charging cycles continuously, there is serious thickening of the SEI film and depletion of active lithium in the battery. After 867 cycles approaching the end-of-life phase, there is more than double ohmic resistance, leading to a change in the system transfer function into:
\begin{equation}
H(s) = 0.0196 + \frac{0.0310}{1 + 10.02s}. \label{eq:aged_tf}
\end{equation}
Due to the increase in electrochemical time constant $\tau$ to $10.02\text{ seconds}$ due to severe interfacial impedance, there is movement of the system pole rightwards to $-0.0998\text{ rad/s}$ from $-0.1285\text{ rad/s}$ on the s-plane, as illustrated in Fig.~\ref{fig:pole_zero_migration}.

\begin{figure}[ht!]
\centering
\includegraphics[width=0.90\textwidth]{pole_zero_migration.png}
\caption{Complex s-plane pole-zero migration tracking severe interfacial impedance growth across the operational lifetime of unseen evaluation test cell \texttt{2017-05-12\_cell\_12}. As ohmic resistance doubles, the system pole migrates rightward toward the instability boundary.}\label{fig:pole_zero_migration}
\end{figure}

\subsection{Closing the Loop: Electrochemical Correlation with Feature Gain}\label{subsec:closing_loop}

The shifting of pole and zero live on the s-plane serves as concrete physical evidence linking our ML model prediction to the actual electrochemical deterioration process. As revealed by the LightGBM feature gain in Section IV, internal resistance ($\text{IR}$) turns out to be the predominant variable with more than $35\%$ tree split contribution. The migration of the system pole $p_1$ to the right in the s-plane is like an electrical heartbeat, serving as mathematical evidence showing why internal resistance is the major reason behind battery deterioration. Due to the $\text{IR}$ doubling caused by SEI layer thickness increase, pole migration takes place towards instability. Terminal voltage sag occurs due to load conditions.

\section{Results and Discussion}\label{sec:results_discussion}

\subsection{Regression Performance across Unseen Test Cohort}\label{subsec:regression_perf}

In order to conduct a thorough assessment of the prediction generalization performance, the predictive power of our gradient boosted LightGBM prediction framework was tested using the 24 commercially available LFP batteries which were kept in strictly hidden fashion using GroupShuffleSplit. Table~\ref{tab:test_performance} presents the empirical regression metrics obtained for the training and test sets respectively. The prediction algorithm achieves a Mean Absolute Error ($\text{MAE}$) of $48.70\text{ cycles}$ and an $R^2$ accuracy of $95.74\%$ on the 100 training battery cells. However, the out-of-sample test battery cells consisting of 4,234 data points achieve an $\text{MAE}$ of $81.35\text{ cycles}$ and an $R^2$ of $78.52\%$.

\begin{table}[ht!]
\centering
\caption{Unseen Test Cohort Validation Performance}\label{tab:test_performance}
\small
\begin{tabular}{@{}l c c c@{}}
\toprule
\textbf{Dataset Split} & \textbf{Cell Count} & \textbf{MAE Error} & \textbf{$R^2$ Accuracy} \\
\midrule
Training Set & 100 Cells & 48.70 cycles & 95.74\% \\[1ex]
Unseen Test Set & 24 Cells & 81.35 cycles & 78.52\% \\
\bottomrule
\end{tabular}
\end{table}

As evident from the parity scatter plot (Fig.~\ref{fig:true_vs_predicted}), predictions are highly correlated along the ideal prediction line throughout the whole range of cycle life. It is normal and scientifically sound to have a performance difference between training ($48.70\text{ MAE}$) and testing ($81.35\text{ MAE}$). This is because, unlike naive random k-fold split, which interpolates consecutive cycles within one cell, we use GroupShuffleSplit to ensure full separation of cells. Inherent physical differences among cells due to different manufacturing, electrolyte composition, and charging protocols make such an out-of-sample $\text{MAE}$ of $\sim 81\text{ cycles}$ a very robust generalization.

\begin{figure}[ht!]
\centering
\includegraphics[width=0.85\textwidth]{01_true_vs_predicted_rul.png}
\caption{Parity scatter plot comparing true observed cycle life against LightGBM predicted remaining useful life across all 24 unseen evaluation test batteries ($\text{MAE} = 81.35\text{ cycles}$, $R^2 = 78.52\%$).}\label{fig:true_vs_predicted}
\end{figure}

\subsection{Baseline Comparison and Empirical Benchmarking}\label{subsec:baseline_comparison}

In order to prove the effectiveness of the LightGBM framework, we compared it to conventional prognostic models, tested on the very same 24 test cells: the linear regression method (the basic Severson model), Random Forest Regressor, and XGBoost Regressor. In order to provide an equal comparison based on computational power and latency, parameters were fine-tuned in Random Forest and XGBoost in order to reach accuracy comparable to that of our model.

\begin{table}[ht!]
\centering
\caption{Empirical Benchmarking of Predictive Models on Unseen Test Cohort}\label{tab:empirical_benchmarking}
\small
\begin{tabular}{@{}l c c c c p{2.8cm}@{}}
\toprule
\textbf{Model} & \textbf{MAE (Cycles)} & \textbf{$R^2$ (\%)} & \textbf{Flash Size} & \textbf{MCU Loop} & \textbf{Automotive ECU Feasibility} \\
\midrule
Linear Regression & 152.90 & 56.12\% & 1.2 KB & 0.22 ms & Deployable but high error ($>150\text{c}$). \\[1ex]
Random Forest & 80.18 & 79.15\% & 28.1 MB & 1.85 ms & \textbf{FAILED}: Exceeds 1 MB Flash limit. \\[1ex]
XGBoost Regressor & 82.53 & 80.25\% & 1.36 MB & 0.68 ms & \textbf{FAILED}: Exceeds 1 MB Flash limit. \\[1ex]
\textbf{LightGBM (Ours)} & \textbf{79.91} & \textbf{81.62\%} & \textbf{809 KB} & \textbf{0.41 ms} & \textbf{OPTIMAL}: Fits Flash, $<1.5\text{ ms}$ loop. \\
\bottomrule
\end{tabular}
\end{table}

According to Table~\ref{tab:empirical_benchmarking} and Fig.~\ref{fig:benchmark_bar}, the Severson linear regression baseline model results in an $\text{MAE}$ of $152.90\text{ cycles}$ ($R^2 = 56.12\%$), proving the incapability of linear feature combinations in predicting the abrupt, non-linear drop in capacity at the end of battery life. The accuracy of models like Random Forest ($\text{MAE} = 80.18$, $R^2 = 79.15\%$) and XGBoost ($\text{MAE} = 82.53$, $R^2 = 80.25\%$) is good, but they do not satisfy crucial automotive hardware requirements. Random Forest needs $28.1\text{ MB}$ of serialized memory space and $1.85\text{ ms}$ MCU traversal time, far beyond the typical $1\text{ MB}$ limit of microcontroller flash memory. XGBoost needs $1.36\text{ MB}$ and $0.68\text{ ms}$ traversal time. On the contrary, our LightGBM framework provides optimal accuracy ($\text{MAE} = 79.91$, $R^2 = 81.62\%$) with just $809\text{ KB}$ serialized flash and $0.41\text{ ms}$ tree traversal time, making it the only feasible algorithm on inexpensive automotive ECUs.

\begin{figure}[ht!]
\centering
\includegraphics[width=0.95\textwidth]{model_comparison_bar_chart.png}
\caption{Empirical benchmarking comparison evaluating predictive accuracy ($\text{MAE}$ and $R^2$) against linear regression, Random Forest, and XGBoost on identical unseen test cells.}\label{fig:benchmark_bar}
\end{figure}

\subsection{Conformal Prediction Coverage Verification}\label{subsec:conformal_verification}

For safety-related applications in electric vehicles, point estimates should be mathematically complemented with uncertainty measures. Through the use of split conformal prediction based on validation residuals, our methodology derives a safety interval around the point estimate of $\pm 122\text{ cycles}$. When checking empirical coverage across all 4,234 test data points, we find that precisely $90.4\%$ of all observations are safely contained within this $\pm 122\text{ cycle}$ window.

\begin{figure}[ht!]
\centering
\includegraphics[width=0.85\textwidth]{03_prediction_errors_histogram.png}
\caption{Distribution of prediction errors ($\text{True} - \text{Predicted cycles}$) showing a zero-centered, symmetrical profile bounded by the $\pm 122\text{ cycle}$ 90\% conformal safety window.}\label{fig:error_histogram}
\end{figure}

Prediction error histogram analysis (Fig.~\ref{fig:error_histogram}) shows a symmetric and zero-based distribution with a small standard deviation. Conservative overestimation happens in early life and is minor, taking place around Cycle 50 when fresh LFP cells have long plateaus at low voltage levels until normal aging wear and tear begins. As the safety margin covers all cases of non-Gaussian skewness, $122\text{ cycles}$ should be subtracted from predicted point estimates to establish a guaranteed worst-case lower bound for maintenance scheduling before failure.

\subsection{Real-Time Dynamic Trajectory Tracking}\label{subsec:trajectory_tracking}

To demonstrate live prognostic tracking across full operational degradation, continuous lifecycle evaluation was performed on unseen test cell \texttt{2017-05-12\_cell\_12}. The LightGBM point prediction dynamically traces the true observed $\text{RUL}$ trajectory from initial deployment down to the 80\% State of Health retirement threshold, bounded continuously by the shaded 90\% conformal prediction interval. Unlike static physics models that diverge during variable load profiles, our data-driven architecture immediately adjusts its slope downward as soon as thermal and internal resistance spikes accelerate aging.

\subsection{Prognostic Maintenance Alert Classification}\label{subsec:alert_classification}

Beyond continuous regression, onboard battery management systems need to classify emergency situations for maintenance alerts whenever a battery cell is nearing its end of life. We define an alert situation when there are less than or equal to 100 life cycles remaining ($\text{RUL} \le 100$). The confusion matrix generated from all 4,234 test points is illustrated in Fig.~\ref{fig:confusion_matrix}.

\begin{figure}[ht!]
\centering
\includegraphics[width=0.75\textwidth]{06_confusion_matrix.png}
\caption{Prognostic maintenance confusion matrix evaluating emergency replacement alert classification ($\text{RUL} \le 100\text{ cycles}$) across 4,234 unseen evaluation checkpoints.}\label{fig:confusion_matrix}
\end{figure}

In the test set, the algorithm classified True Healthy states ($>100\text{ cycles}$) 3,600 times (True Negatives, $\text{TN}$) with just 72 False Positives ($\text{FP}$). The model correctly identified critical aging states ($\le 100\text{ cycles}$) 498 times (True Positives, $\text{TP}$), failing to do so only 64 times (False Negatives, $\text{FN}$). Using these empirical confusion matrix counts across all 4,234 evaluation points, we calculate the formal classification performance metrics along with their exact numerical results as follows:

\begin{align}
\text{Accuracy} &= \frac{\text{TP} + \text{TN}}{\text{TP} + \text{TN} + \text{FP} + \text{FN}} = \frac{498 + 3600}{498 + 3600 + 72 + 64} = \frac{4098}{4234} = 96.79\%, \label{eq:accuracy} \\[1ex]
\text{Precision} &= \frac{\text{TP}}{\text{TP} + \text{FP}} = \frac{498}{498 + 72} = \frac{498}{570} = 87.37\%, \label{eq:precision} \\[1ex]
\text{Recall} &= \frac{\text{TP}}{\text{TP} + \text{FN}} = \frac{498}{498 + 64} = \frac{498}{562} = 88.61\%, \label{eq:recall} \\[1ex]
F_1 &= 2 \cdot \frac{\text{Precision} \cdot \text{Recall}}{\text{Precision} + \text{Recall}} = 2 \cdot \frac{0.8737 \cdot 0.8861}{0.8737 + 0.8861} = 87.98\%. \label{eq:f1_score}
\end{align}
An $F_1$-score of $87.98\%$ is evidence that the presented framework indeed offers an extremely high level of decision-making support by reducing the chances of either replacing batteries prematurely or facing unexpected battery failure on the road.

\subsection{Computational Feasibility and Microcontroller Benchmarking}\label{subsec:mcu_benchmarking}

The ability to run in real-time on automotive hardware is one of the key contributions made by this research. The timing information for running on a simulated 100 MHz ARM Cortex microcontroller is shown in Table~\ref{tab:mcu_timing}. Sensor value gathering from CAN bus registers takes $0.12\text{ ms}$. Feature extraction within the rolling window (calculation of $\text{d}Q\_\log\_\text{var}$ and capacity fade for $W=10$ windows) takes $0.42\text{ ms}$. LightGBM leaf-wise tree evaluation across 300 decision trees requires $0.41\text{ ms}$ via simple integer threshold operations.

\begin{table}[ht!]
\centering
\caption{Hardware Microcontroller Benchmarking \& Timing Breakdown}\label{tab:mcu_timing}
\small
\begin{tabular}{@{}l c c p{3.5cm}@{}}
\toprule
\textbf{Execution Step} & \textbf{Operations Performed} & \textbf{Measured Latency} & \textbf{Hardware Microchip Feasibility} \\
\midrule
1. Sensor Data Acquisition & Reading V, I, T buffers & 0.12 ms & Direct CAN bus register read \\[1ex]
2. Feature Extraction & Rolling $\text{d}Q\_\log\_\text{var}$ ($W=10$) & 0.42 ms & Simple integer rolling variance \\[1ex]
3. Tree Evaluation & Traversing 300 trees & 0.41 ms & Integer IF/ELSE threshold logic \\[1ex]
\midrule
\textbf{TOTAL INFERENCE} & \textbf{Full End-to-End Prediction} & \textbf{0.95 ms} & \textbf{Leaves $>99.9\%$ CPU free ($<1.5\text{ ms}$ budget)} \\
\bottomrule
\end{tabular}
\end{table}

Total end-to-end inference is completed in exactly $0.95\text{ ms}$, safely under the $1.5\text{ ms}$ real-time control system budget with $>99.9\%$ of CPU cycles available for vehicle control operations. Also, when compared to deep learning, the decision tree architecture enjoys superior speed of operation because recurrent LSTM neural nets take approximately $45\text{ ms}$ for each inference due to floating-point matrix multiplication while our LightGBM algorithm executes $47\times$ faster. Moreover, thanks to the small compiled footprint of approximately $480\text{ KB}$, the model fits effortlessly into standard automotive microcontrollers.

\section{Conclusion}\label{sec13}

Real-time and accurate prognosis of Lithium Iron Phosphate (LFP) batteries is still a challenging task in the field of electric vehicles because of their characteristic voltage plateau that is indicative of internal degradation of active material. In this paper, a physics-based machine learning framework, using gradient-boosted decision trees (LightGBM), was introduced specifically designed for Battery Management Systems (BMS) in automobiles. Our approach avoids expensive deep learning algorithms and uses voltage measurement.

The comprehensive analysis using 124 commercialized APR18650M1A cylindrical LFP cells (with a total of 22,474 operational evaluation points based on 72 different multi-step fast charging policies) clearly revealed the superior performance of our approach. By using a leak-proof GroupShuffleSplit protocol for cross-validation at the physical cell level, we found that the LightGBM framework provided out-of-sample Mean Absolute Error ($\text{MAE}$) of $81.35\text{ cycles}$ and an $R^2$ accuracy of $78.52\%$ with 24 strictly unseen batteries. Moreover, empirical benchmarking showed the superiority of our model compared to linear regression ($\text{MAE}$ of $152.90$) and comparable performance compared to Random Forest ($\text{MAE}$ of $80.18$) and XGBoost ($\text{MAE}$ of $82.53$).

Importantly, profiling on simulated automotive ARM Cortex microcontrollers showed that inference using the LightGBM model takes precisely $0.95\text{ milliseconds}$---$47$ times faster than recurrent LSTM models (around $45\text{ ms}$)---using just $809\text{ KB}$ of compiled flash memory space. This enables the algorithm to be seamlessly deployed in standard $1\text{ MB}$ automotive Electronic Control Units (ECUs) without affecting real-time operations in the vehicle control processes. Furthermore, split conformal prediction was able to establish safe uncertainty bounds of $\pm 122$ clock cycles with a $90.4\%$ coverage rate.

Lastly, the mathematically proven importance rankings of our data-driven approach, which found internal resistance ($\text{IR}$) to be the most prominent split feature ($> 35\%$ contribution), was validated through first order Equivalent Circuit Model ($\text{ECM}$) transfer function analysis. Real-time monitoring of the complex s-plane dynamics showed that as cell interfacial resistance doubles during aging, the system pole moves rightwards from $-0.1285\text{ rad/s}$ towards instability, which proved the data-driven model's ability to capture the underlying physics of the problem. Future work will explore how to leverage the physics-informed architecture of the LightGBM model to enable adaptive transfer learning for different chemistries in silicon-anode and sodium-ion cells as well as cloud-connected OTA calibrations in commercial EV fleets.

\nocite{*}
\bibliography{sn-bibliography}% common bib file

\end{document}
